% ============================================================
\setpartpreamble{\partintro
  {Limits of the run-time hooks, further extension points, operating in production and the
  way from the plugin into the application.}
  {
    \item Chapter~\ref{ch:bulk-operations-and-the-limits}: Bulk operations and the limits of
    interceptors
    \item Chapter~\ref{ch:value-generation}: Value generation
    \item Chapter~\ref{ch:the-update-pipeline}: The update pipeline
    \item Chapter~\ref{ch:complex-types-and-json}: Complex types and JSON
    \item Chapter~\ref{ch:models-per-runtime-context}: Models per runtime context
    \item Chapter~\ref{ch:advanced-migrations}: Advanced migrations
    \item Chapter~\ref{ch:row-level-security-in-production}: Row Level Security in
    production
    \item Chapter~\ref{ch:tenant-safe-relationships}: Tenant-safe relationships
    \item Chapter~\ref{ch:retrofitting-and-operating}: Retrofitting and operating
    multi-tenancy
    \item Chapter~\ref{ch:plugin-performance}: Plugin performance
    \item Chapter~\ref{ch:using-the-plugin-in-asp}: Using the plugin in ASP.NET Core
    \item Chapter~\ref{ch:ci-and-api-stability}: CI and API stability
  }}
\part{Advanced topics}\label{part:VIII}

% ------------------------------------------------------------
\chapter{Bulk operations and the limits of interceptors}\label{ch:bulk-operations-and-the-limits}

A SaveChanges interceptor only sees what goes through \code{Save\allowbreak{}Changes}.
\code{Execute\allowbreak{}Update} and \code{Execute\allowbreak{}Delete} write directly to
the database, without change tracker and without SaveChanges. The same applies to raw SQL. A
plugin that enforces its guarantees only through SaveChanges has a gap here:

\begin{datatable}{H{1.28}H{0.88}H{0.88}H{0.96}}
  \textbf{Path} & \textbf{SaveChanges interceptor} & \textbf{Connection interceptor} & \textbf{RLS} \\
  \tablehead
  \code{Save\allowbreak{}Changes} & yes & yes & yes \\
  \tablerule
  \code{Execute\allowbreak{}Update}, \code{Execute\allowbreak{}Delete} & \textbf{no} & yes &
  yes \\
  \tablerule
  \code{From\allowbreak{}Sql}, \code{Sql\allowbreak{}Query}, \code{Execute\allowbreak{}Sql}
  & \textbf{no} & yes & yes \\
  \tablerule
  ADO.NET or psql as the app role & \textbf{no} & \textbf{no} & yes (without tenant: no
  rows) \\
\end{datatable}

In the sample, RLS catches every one of these paths because the database enforces the rules
itself. An integration test proves it by deliberately bypassing the interceptor:

% source: tests/Acme.EntityFrameworkCore.Tenancy.IntegrationTests/AdvancedTests.cs
\begin{sourcecode}[tests/\allowbreak{}.\allowbreak{}.\allowbreak{}.\allowbreak{}Integration\allowbreak{}Tests/\allowbreak{}Advanced\allowbreak{}Tests.\allowbreak{}cs (\allowbreak{}excerpt)]{csharp}
[Fact]
public async Task ExecuteUpdate_bypasses_interceptor_but_not_rls()
{
    // ... create an order for "tenant" ...
    await using var attacker = ContextFor(tenant);

    // No SaveChanges -> the SaveChanges interceptor does NOT run.
    // The policy's WITH CHECK still prevents the reassignment.
    var ex = await Assert.ThrowsAsync<PostgresException>(() =>
        attacker.Orders.ExecuteUpdateAsync(
            s => s.SetProperty(o => o.TenantId, TenantId.New())));

    Assert.Equal(PostgresErrorCodes.InsufficientPrivilege, ex.SqlState);
}
\end{sourcecode}

\begin{important}[What a plugin must document]
  Plugins without database-level protection (soft delete, audit fields, tenant filtering via
  query filter only) must document that bulk operations and raw SQL bypass their rules. To
  prevent it, an \code{IQuery\allowbreak{}Expression\allowbreak{}Interceptor} can be used:
  \code{Execute\allowbreak{}Update} and \code{Execute\allowbreak{}Delete} go through the
  query pipeline, the interceptor sees their expression tree and can, for example, reject a
  \code{Set\allowbreak{}Property} on protected columns. Raw SQL remains open regardless.
\end{important}

\section{Example: protecting the tenant column in bulk updates}

The sample implements exactly this interceptor, as an option:
\code{Use\allowbreak{}Acme\allowbreak{}Tenancy(\allowbreak{}t => t.\allowbreak{}Protect\allowbreak{}Bulk\allowbreak{}Operations(\allowbreak{}))}
for PostgreSQL, \code{sql.\allowbreak{}Protect\allowbreak{}Bulk\allowbreak{}Operations(\allowbreak{})} for SQL
Server. EF Core~10 no longer takes the setters of \code{Execute\allowbreak{}Update} as an
expression tree but as an \code{Action<\allowbreak{}Update\allowbreak{}Setters\allowbreak{}Builder<\allowbreak{}T>>}.
In the tree the interceptor sees, they then appear as an array of tuples:
\code{Execute\allowbreak{}Update(\allowbreak{}source,\allowbreak{} new[] \{ new Tuple(\allowbreak{}e => e.\allowbreak{}Tenant\allowbreak{}Id,\allowbreak{} value),\allowbreak{} .\allowbreak{}.\allowbreak{}. \})}.
The first element of each tuple is the property lambda, and that is what the visitor checks:

% source: src/Acme.EntityFrameworkCore.Tenancy/Interceptors/TenantQueryExpressionInterceptor.cs
\begin{sourcecode}[Interceptors/\allowbreak{}Tenant\allowbreak{}Query\allowbreak{}Expression\allowbreak{}Interceptor.\allowbreak{}cs (\allowbreak{}core package,\allowbreak{} shortened)]{csharp}
protected override Expression VisitMethodCall(MethodCallExpression node)
{
    if (node.Method.DeclaringType == typeof(EntityFrameworkQueryableExtensions)
        && node.Method.Name is nameof(EntityFrameworkQueryableExtensions.ExecuteUpdate)
            or nameof(EntityFrameworkQueryableExtensions.ExecuteUpdateAsync)
        && node.Arguments.Count > 1
        && node.Arguments[1] is NewArrayExpression setters)
    {
        foreach (var setter in setters.Expressions)
        {
            if (setter is NewExpression { Arguments: [var selector, ..] }
                && Unquote(selector) is LambdaExpression lambda)
                Check(lambda);
        }
    }

    return base.VisitMethodCall(node);
}

private void Check(LambdaExpression selector)
{
    var parameter = selector.Parameters[0];
    var name = selector.Body switch
    {
        // e => e.TenantId
        MemberExpression member when member.Expression == parameter => member.Member.Name,
        // e => EF.Property<TenantId>(e, "TenantId"), also for shadow properties
        MethodCallExpression
        {
            Method.Name: nameof(EF.Property),
            Arguments: [var target, ConstantExpression { Value: string property }]
        } when target == parameter => property,
        _ => null,
    };
    // ...
}
\end{sourcecode}

If the property being set is the tenant property according to the runtime annotation, the
interceptor throws a \code{Tenancy\allowbreak{}Exception}, before any connection is opened. It runs
once per query shape, before translation and the query cache; later executions cost nothing.
Other columns and ordinary queries pass unchanged. RLS remains the actual boundary: the
interceptor only stops reassigning rows early and with a clear message; it does not check
\code{Execute\allowbreak{}Delete} or raw SQL.

The registration holds a trap. \code{IQuery\allowbreak{}Expression\allowbreak{}Interceptor}, like
\code{IMaterialization\allowbreak{}Interceptor}, derives from
\code{ISingleton\allowbreak{}Interceptor}. Registered under that type in the internal provider,
however, EF never calls it: EF takes query expression interceptors from the
\code{IInterceptor} registrations and the options, materialization interceptors also from the
\code{ISingleton\allowbreak{}Interceptor} registrations. The plugin therefore registers it as
\code{IInterceptor}; stateless as it is, an instance per context does no harm.
\code{Bulk\allowbreak{}Operation\allowbreak{}Tests} records both registrations, the integration
tests show an allowed and a rejected update against PostgreSQL.

% ------------------------------------------------------------
\chapter{Value generation}\label{ch:value-generation}

EF generates key values either in the database (identity, sequences) or client-side via a
\code{Value\allowbreak{}Generator}. For \code{Guid} keys EF does this by itself, for custom
types like \code{Tenant\allowbreak{}Id} it does not. A plugin closes this gap with a
generator, a factory and a convention:

% source: src/Acme.EntityFrameworkCore.Tenancy/ValueGeneration/TenantIdValueGenerator.cs
\begin{sourcecode}[Value\allowbreak{}Generation/\allowbreak{}Tenant\allowbreak{}Id\allowbreak{}Value\allowbreak{}Generator.\allowbreak{}cs (\allowbreak{}core package)]{csharp}
public class TenantIdValueGenerator : ValueGenerator<TenantId>
{
    // Final values, no temporary placeholders until saving
    public override bool GeneratesTemporaryValues => false;

    public override TenantId Next(EntityEntry entry) => TenantId.New();   // UUIDv7
}
\end{sourcecode}

% source: src/Acme.EntityFrameworkCore.Tenancy/ValueGeneration/TenantIdValueGeneratorFactory.cs
\begin{sourcecode}[Value\allowbreak{}Generation/\allowbreak{}Tenant\allowbreak{}Id\allowbreak{}Value\allowbreak{}Generator\allowbreak{}Factory.\allowbreak{}cs (\allowbreak{}core package)]{csharp}
public class TenantIdValueGeneratorFactory : ValueGeneratorFactory
{
    // Stateless generator: EF may cache it and share it between contexts
    public override ValueGenerator Create(IProperty property, ITypeBase typeBase)
        => new TenantIdValueGenerator();
}
\end{sourcecode}

% source: src/Acme.EntityFrameworkCore.Tenancy/Conventions/TenantIdKeyValueGenerationConvention.cs
\begin{sourcecode}[Conventions/\allowbreak{}Tenant\allowbreak{}Id\allowbreak{}Key\allowbreak{}Value\allowbreak{}Generation\allowbreak{}Convention.\allowbreak{}cs (\allowbreak{}core package)]{csharp}
public virtual void ProcessModelFinalizing(
    IConventionModelBuilder modelBuilder,
    IConventionContext<IConventionModelBuilder> context)
{
    foreach (var entityType in modelBuilder.Metadata.GetEntityTypes())
    {
        if (entityType.BaseType is not null
            || entityType.FindPrimaryKey() is not { Properties: [var property] }
            || property.ClrType != typeof(TenantId))
            continue;

        // Convention source: an explicit ValueGeneratedNever() wins (-> null)
        if (property.Builder.ValueGenerated(ValueGenerated.OnAdd) is null)
            continue;

        property.Builder.HasValueGeneratorFactory(typeof(TenantIdValueGeneratorFactory));
    }
}
\end{sourcecode}

\begin{important}[A factory type, not a generator]
  A value generator can be configured in several ways:
  \code{Has\allowbreak{}Value\allowbreak{}Generator} with a type or a factory function, or
  \code{Has\allowbreak{}Value\allowbreak{}Generator\allowbreak{}Factory} with a factory
  type. EF Core 10's compiled model generator only accepts the last one and otherwise aborts
  with “The property ... has a value generator configured. Use 'HasValueGeneratorFactory' to
  configure the value generator factory type.” Plugins should therefore use
  \code{Has\allowbreak{}Value\allowbreak{}Generator\allowbreak{}Factory(\allowbreak{}Type)}
  and give the factory a public parameterless constructor.
\end{important}

The convention lives in the core package because it has nothing to do with PostgreSQL; the
provider package plugs it in via its convention set plugin. The compiled model test checks
that the factory appears in the generated code.

% ------------------------------------------------------------
\chapter{The update pipeline}\label{ch:the-update-pipeline}

Chapter~\ref{ch:the-query-pipeline} describes the way from LINQ to SQL. Writing has its own
pipeline, which plugins extend for bulk operations, database-side soft deletes or returned
values:

\begin{flowdiagram}
  \flowstep{ChangeTracker}{Determines added, modified and deleted entities}
  \flowstep{IModificationCommandFactory}{One ModificationCommand per row and table}
  \flowstep{ICommandBatchPreparer}{Sorts by dependencies (foreign keys) and forms batches}
  \flowstep{IModificationCommandBatchFactory}{Creates the batches (size, provider
  specifics)}
  \flowstep{IUpdateSqlGenerator}{Writes INSERT, UPDATE, DELETE, possibly with RETURNING}
  \flowstep{Execution}{Results (identity, concurrency tokens) back into the entities}
\end{flowdiagram}

\begin{datatable}{H{1.64}H{0.68}H{0.68}}
  \textbf{Service} & \textbf{Lifetime} & \textbf{Kind} \\
  \tablehead
  \code{IModification\allowbreak{}Command\allowbreak{}Factory} & Singleton & single \\
  \tablerule
  \code{ICommand\allowbreak{}Batch\allowbreak{}Preparer} & Scoped & single \\
  \tablerule
  \code{IModification\allowbreak{}Command\allowbreak{}Batch\allowbreak{}Factory} & Scoped &
  single \\
  \tablerule
  \code{IUpdate\allowbreak{}Sql\allowbreak{}Generator} & Singleton & single \\
\end{datatable}

All of these are single provider services, with the consequences from
chapter~\ref{ch:registering-services}. Before replacing them, it pays to check the
alternatives: a SaveChanges interceptor can rewrite states (delete becomes modify for soft
delete), a command interceptor can add SQL, and for bulk changes there is
\code{Execute\allowbreak{}Update}. A custom
\code{IUpdate\allowbreak{}Sql\allowbreak{}Generator} is only justified when the SQL itself
needs a different shape, for example \code{INSERT .\allowbreak{}.\allowbreak{}.\allowbreak{}
ON CONFLICT} for upserts.

% ------------------------------------------------------------
\chapter{Complex types and JSON}\label{ch:complex-types-and-json}

\section{Complex types}

Complex types (since EF Core 8) are value types without their own identity whose properties
end up in the owner's table. For plugins this means: conventions iterating
\code{entity\allowbreak{}Type.\allowbreak{}Get\allowbreak{}Properties(\allowbreak{})} do
\textbf{not} see their properties. They hang off \code{IComplex\allowbreak{}Property} and
must be iterated separately:

\begin{sourcecode}[Pattern for conventions (\allowbreak{}sketch)]{csharp}
foreach (var complexProperty in entityType.GetComplexProperties())
{
    foreach (var property in complexProperty.ComplexType.GetProperties())
    {
        // Column names are prefixed (e.g. "Address_City"); always determine them via
        // property.GetColumnName(table), never assemble them yourself
    }
}
\end{sourcecode}

Fluent API methods for properties should additionally be offered on
\code{Complex\allowbreak{}Type\allowbreak{}Property\allowbreak{}Builder} where they make
sense there. The sample plugin does not need this: isolation and value generation only
concern entities.

\section{JSON columns}

Owned types and (since EF Core 10) complex types can be stored as a single JSON column with
\code{To\allowbreak{}Json(\allowbreak{})}. Values in it are not read and written through
converters but through the type mapping's
\code{Json\allowbreak{}Value\allowbreak{}Reader\allowbreak{}Writer}
(chapter~\ref{ch:type-mappings-and-value-converters}). Without it, model building fails. The
sample's unit and integration tests check both sides:

% source: tests/Acme.EntityFrameworkCore.Tenancy.Tests/JsonTests.cs
\begin{sourcecode}[tests/\allowbreak{}.\allowbreak{}.\allowbreak{}.\allowbreak{}/\allowbreak{}Json\allowbreak{}Tests.\allowbreak{}cs]{csharp}
using var ctx = Contexts.Create(model: m =>
    m.Entity<Shipment>().OwnsOne(s => s.Details, d => d.ToJson()));

var property = ctx.DesignModel()
    .FindEntityType(typeof(ShipmentDetails))!
    .FindProperty(nameof(ShipmentDetails.ApprovedBy))!;

Assert.Same(TenantIdJsonReaderWriter.Instance,
    property.GetTypeMapping().JsonValueReaderWriter);
\end{sourcecode}

JSON outside the database is a separate matter. System.Text.Json does not know EF's
\code{Json\allowbreak{}Value\allowbreak{}Reader\allowbreak{}Writer} and would write
\code{TenantId} as an object with \code{Value} and \code{IsEmpty}. A
\code{[Json\allowbreak{}Converter(\allowbreak{}typeof(\allowbreak{}Tenant\allowbreak{}Id\allowbreak{}Json\allowbreak{}Converter))]}
on the type in the core package turns it into a plain GUID, also as a dictionary key, without a
web API having to register anything; invalid values lead to a
\code{Json\allowbreak{}Exception}. Both ways need a building block of their own, because EF and
System.Text.Json serialize independently of each other
(\code{Tenant\allowbreak{}Id\allowbreak{}Json\allowbreak{}Converter\allowbreak{}Tests}).

\begin{note}
  PostgreSQL arrays (\code{List<\allowbreak{}Tenant\allowbreak{}Id>} as \code{uuid[]}) are
  not JSON: Npgsql maps primitive collections natively as arrays. Whether a provider
  supports collections with converted element types is provider-specific and deserves its
  own test before you document it.
\end{note}

% ------------------------------------------------------------
\chapter{Models per runtime context}\label{ch:models-per-runtime-context}

EF builds the model once per context type and caches it
(chapter~\ref{ch:ef-core-s-architecture-from}). If the model should depend on run-time data,
such as one schema per tenant, the cache key must contain that data. That is what
\code{IModel\allowbreak{}Cache\allowbreak{}Key\allowbreak{}Factory} is for:

\begin{sourcecode}[tests/\allowbreak{}.\allowbreak{}.\allowbreak{}.\allowbreak{}/\allowbreak{}Schema\allowbreak{}Per\allowbreak{}Tenant\allowbreak{}Tests.\allowbreak{}cs (\allowbreak{}shortened, simplified)]{csharp}
public sealed class SchemaModelCacheKeyFactory : IModelCacheKeyFactory
{
    public object Create(DbContext context, bool designTime)
        => context is ISchemaScoped scoped
            ? (context.GetType(), scoped.Schema, designTime)
            : (object)(context.GetType(), designTime);
}

public sealed class SchemaContext : DbContext, ISchemaScoped
{
    // Constructor: (DbContextOptions<SchemaContext> options, string schema)
    public string Schema { get; }

    protected override void OnModelCreating(ModelBuilder modelBuilder)
        => modelBuilder.HasDefaultSchema(Schema);
}

// Registration
options.UseNpgsql(connectionString)
       .ReplaceService<IModelCacheKeyFactory, SchemaModelCacheKeyFactory>();
\end{sourcecode}

The test \code{Schema\allowbreak{}Per\allowbreak{}Tenant\allowbreak{}Tests} checks both
sides of the cache key: two schemas produce two models with their own default schema, and
the same schema reuses the cached model instead of building it again.

This is the basis for “schema per tenant”, an alternative to RLS. The three usual approaches
compared:

\begin{datatable}{H{0.83}H{1.02}H{1.15}}
  \textbf{Approach} & \textbf{Strengths} & \textbf{Weaknesses} \\
  \tablehead
  Shared tables + RLS (sample) & One model, one schema, simple migrations; the database
  enforces separation & All tenants share tables and indexes; misconfigured roles undermine
  separation \\
  \tablerule
  Schema per tenant & Strong separation, separate backups per schema possible & One model
  per tenant in the cache (memory!), migrations for every schema \\
  \tablerule
  Database per tenant & Maximum separation, independent scaling & One model suffices, but
  connections, migrations and operations per database \\
\end{datatable}

\begin{important}[Model cache and memory]
  Every key value produces its own model in memory. With hundreds of tenants and large
  models this becomes expensive. In addition, values that change the model must never come
  from unchecked input.
\end{important}

% ------------------------------------------------------------
\chapter{Advanced migrations}\label{ch:advanced-migrations}

\section{Order of operations}

Generated and hand-written operations run in the order they appear in
\code{Up(\allowbreak{})}. This is a real source of errors. In the sample, the policy
\code{support\_\allowbreak{}read} applies \code{TO shop\_\allowbreak{}support}. EF generates
it directly after \code{CREATE TABLE}. If the role is only created at the end of
\code{Up(\allowbreak{})}, the migration fails.

\begin{sourcecode}[Migrations/\allowbreak{}20261021\_\allowbreak{}Initial.\allowbreak{}cs]{csharp}
protected override void Up(MigrationBuilder migrationBuilder)
{
    // 1. Whatever policies reference: BEFORE the generated operations
    migrationBuilder.CreateApplicationRole("shop_support");

    // 2. Generated by EF: tables, RLS, policies (TO shop_support)
    migrationBuilder.CreateTable(name: "Customers", /* ... */);

    // 3. Whatever requires the tables: AFTER the generated operations
    migrationBuilder.CreateApplicationRole("shop_app", canLogin: true);
    migrationBuilder.GrantTablePrivileges("shop_app");     // applies to existing tables
}
\end{sourcecode}

\section{Idempotent scripts}

\code{dotnet ef migrations script -{}-idempotent} produces a script that can run against any
database state. For this, Npgsql wraps every operation in a \code{DO \$EF\$
.\allowbreak{}.\allowbreak{}.\allowbreak{} \$EF\$} block that checks the history table.
Custom operations that produce a \code{DO \$\$} block themselves end up inside it. This
works because the dollar quotes use different tags. Assumptions like this belong in a test:

% source: tests/Acme.EntityFrameworkCore.Tenancy.IntegrationTests/AdvancedTests.cs
\begin{sourcecode}[tests/\allowbreak{}.\allowbreak{}.\allowbreak{}.\allowbreak{}Integration\allowbreak{}Tests/\allowbreak{}Advanced\allowbreak{}Tests.\allowbreak{}cs (\allowbreak{}excerpt)]{csharp}
var script = db.GetService<IMigrator>()
    .GenerateScript(options: MigrationsSqlGenerationOptions.Idempotent);

await db.Database.ExecuteSqlRawAsync(script);   // nested DO valid?
await db.Database.ExecuteSqlRawAsync(script);   // second run is skipped
\end{sourcecode}

\section{Operations outside of transactions}

Some commands must not run inside a transaction, in PostgreSQL for example \code{CREATE
INDEX CONCURRENTLY}. Custom operations with such SQL end their command with
\code{End\allowbreak{}Statement(\allowbreak{}builder,\allowbreak{}
suppress\allowbreak{}Transaction: true)}. EF then executes it outside the migration
transaction. The price: if the migration fails afterwards, this command is not rolled back;
it should therefore be idempotent.

\section{Migration bundles}

\code{dotnet ef migrations bundle} produces an executable that applies migrations without
the SDK. It contains the startup project including the plugin assemblies; custom operations
and the plugin's SQL generator therefore take effect there too. It is important that the
design-time factory or host uses the same plugin options as the application, otherwise SQL
for a different model is produced.

% ------------------------------------------------------------
\chapter{Row Level Security in production}\label{ch:row-level-security-in-production}

RLS is only as strong as its operating conditions. These pitfalls silently render the
policies ineffective or unusable:

\begin{datatable}{H{0.93}H{0.93}H{1.14}}
  \textbf{Pitfall} & \textbf{Consequence} & \textbf{Remedy} \\
  \tablehead
  App runs as table owner & Policies are ignored & \code{FORCE ROW LEVEL SECURITY} (plugin
  default) and separate roles \\
  \tablerule
  Superuser or \code{BYPASSRLS} & RLS is always bypassed, even with FORCE & App role with
  \code{NOSUPERUSER NOBYPASSRLS} (custom operation, chapter~\ref{ch:custom-migration-operations}) \\
  \tablerule
  Setting not set & Without \code{nullif}: error; with it: empty results & The interceptor
  sets the value on every open \\
  \tablerule
  No index on the tenant column & Every query scans the whole table & Index convention
  (chapter~\ref{ch:conventions}) \\
  \tablerule
  Transaction pooling (PgBouncer) & Session settings reach other clients or get lost & Set
  the setting transaction-locally (see below) \\
  \tablerule
  Connection opened outside EF & The session keeps an old tenant or none & Let EF open
  connections, or \code{Apply\allowbreak{}Tenant\allowbreak{}Session(\allowbreak{})}
  (chapter~\ref{ch:interceptors}) \\
  \tablerule
  “Bypass” switch for the app role & Anyone who can run SQL as the app role can flip it &
  Separate roles for cross-tenant access (see below) \\
\end{datatable}

\section{The threat model}\label{sec:threat-model}

RLS is a layer of protection with clear limits. Whoever uses it should know what it protects
against, what it relies on and where it ends:

\begin{datatable}{H{0.45}H{1.55}}
  \textbf{Aspect} & \textbf{In the sample} \\
  \tablehead
  Protected & The rows of tenant-owned tables in the supported mappings
  (page~\pageref{sec:supported-table-mappings}). Not protected are shared master data, the
  schema and, without composite keys, the existence of foreign values
  (chapter~\ref{ch:tenant-safe-relationships}). \\
  \tablerule
  Trusted & The database server and its roles; the migration role as table owner; the code of
  the application and the plugin, including the interceptors; the
  \code{ITenant\allowbreak{}Provider} and the source it reads the tenant from. \\
  \tablerule
  Assumptions & The app role is neither superuser nor owner and has no \code{BYPASSRLS}. The
  tenant comes from an authenticated claim, not from a header or parameter the client chooses
  freely. EF opens the connections, and there is no transaction pooling. \\
  \tablerule
  Limit & The app role may set the session setting itself; it has to, so that the interceptor
  can set it. Whoever runs arbitrary SQL under this role switches the tenant with it. \\
\end{datatable}

The limit is not a flaw of the plugin but a property of the session model.
\code{Threat\allowbreak{}Model\allowbreak{}Tests} shows it against PostgreSQL: one
\code{SELECT set\_\allowbreak{}config(\allowbreak{}'app.\allowbreak{}tenant\_\allowbreak{}id',
\allowbreak{} .\allowbreak{}.\allowbreak{}.\allowbreak{})} under the app role is enough, and the
query after it returns the other tenant's rows.
\code{Sql\allowbreak{}Server\allowbreak{}Threat\allowbreak{}Model\allowbreak{}Tests} shows the
same with \code{sp\_\allowbreak{}set\_\allowbreak{}session\_\allowbreak{}context}. This defines
what RLS protects against here and what it does not:

\begin{itemize}
  \item \textbf{Protected:} forgotten filters, faulty LINQ queries, the application's raw SQL
    with a wrong \code{WHERE}, bulk operations, reports that accidentally run through the app
    role. In all these cases the plugin sets the tenant, and the code does not change it.
  \item \textbf{Not protected:} SQL injection and compromised application code. Both can change
    the setting, even within a single statement, since \code{set\_\allowbreak{}config} is a
    function. What helps against them are parameters instead of concatenated strings (EF
    parameterizes LINQ and \code{From\allowbreak{}Sql}/\code{Sql\allowbreak{}Query} with
    interpolated strings; \code{From\allowbreak{}Sql\allowbreak{}Raw} with concatenated values
    does not) and privileges that give the app role only the tables it needs.
  \item \textbf{Outside the database:} a forged or wrongly assigned claim grants access to
    exactly that tenant. That is the job of authentication; the plugin takes the tenant from
    the \code{ITenant\allowbreak{}Provider}.
\end{itemize}

The sample deliberately does not implement possible hardenings. SQL Server can freeze the key
until the end of the connection with \code{sp\_\allowbreak{}set\_\allowbreak{}session\_\allowbreak{}context
.\allowbreak{}.\allowbreak{}.\allowbreak{}, @read\_only = 1}
(\href{https://learn.microsoft.com/en-us/sql/relational-databases/security/row-level-security}{SQL Server: Row-Level Security});
the test above shows that a later change then fails. The price: \code{Apply\allowbreak{}Tenant\allowbreak{}Session},
too, can no longer change the tenant of an open connection. PostgreSQL has no read-only setting;
there remain a role per tenant with \code{SET ROLE} or a signed tenant token checked by a
\code{SECURITY DEFINER} function. Both move complexity into the database and are only named as a
direction here. Cross-tenant access still belongs in roles of its own (section below).

\section{Transaction pooling}

PgBouncer in “transaction” mode shares a server connection between clients per transaction.
A session value set with
\code{set\_\allowbreak{}config(\allowbreak{}.\allowbreak{}.\allowbreak{}.\allowbreak{},\allowbreak{}
false)} can then reach another client, a security hole. The safe variant: set the setting
\emph{transaction-locally} with
\code{set\_\allowbreak{}config(\allowbreak{}.\allowbreak{}.\allowbreak{}.\allowbreak{},\allowbreak{}
true)} and run all database work in an explicit transaction. Technically, the interceptor
moves from \code{IDb\allowbreak{}Connection\allowbreak{}Interceptor} to
\code{IDb\allowbreak{}Transaction\allowbreak{}Interceptor}
(\code{Transaction\allowbreak{}Started}). Because queries without a transaction then have no
tenant at all, the application must adhere to this consistently. The sample stays with the
simpler session model and documents that it requires session or Npgsql pooling.

\section{Policy performance}

PostgreSQL inserts policies into every query like additional \code{WHERE} conditions.
\code{current\_\allowbreak{}setting} is \code{STABLE}. That guarantees the same result for
all rows within one statement and \emph{allows} the optimizer to combine several calls into
one (\href{https://www.postgresql.org/docs/17/xfunc-volatility.html}{PostgreSQL: Function Volatility Categories}); it does not guarantee a fixed
number of calls. How often the function actually runs depends on the plan. If the comparison
is part of an \code{Index Cond}, PostgreSQL evaluates the comparison value once and reads only
the tenant's rows. If it is a \code{Filter} over a seq scan, it may be evaluated per row, and
then the scan weighs more than the function. Expensive policies come from subqueries (for
example permissions from another table). In addition, PostgreSQL holds back conditions of the
query whose functions are not \code{LEAKPROOF} behind the policy conditions, which can prevent
optimizations (\href{https://www.postgresql.org/docs/17/ddl-rowsecurity.html}{PostgreSQL: Row Security Policies}).

Only the plan under real conditions is therefore meaningful: \code{EXPLAIN
(\allowbreak{}ANALYZE,\allowbreak{} BUFFERS)} as the app role with a tenant set, with a
realistic amount of data and current statistics. Measured as owner or superuser, the policies
are missing from the plan, because superusers and roles with \code{BYPASSRLS} bypass them, and
so does the owner without \code{FORCE}. The overhead shows in comparing the same query with
and without the policy, not in the number of function calls.

\section{Legitimate cross-tenant access}

Sooner or later someone needs data across tenants: reports for the operator, maintenance
jobs, support. The tempting shortcut is a switch in the application, for example an
additional session setting \code{app.\allowbreak{}bypass} that the policy honors. That is
not a security boundary: whatever the application role can set, an injected SQL statement
can set too. Separation must come from database roles instead:

\begin{datatable}{H{0.75}H{1.25}H{1.00}}
  \textbf{Need} & \textbf{Solution} & \textbf{Why it is safe} \\
  \tablehead
  Reports, analytics & Separate read-only role with \code{BYPASSRLS}, only \code{SELECT},
  its own connection string and a context configuration without tenant & The app role cannot
  become this role; its credentials live elsewhere \\
  \tablerule
  Support reads selected data & Additional permissive policy for a support role:
  \code{Has\allowbreak{}Policy(\allowbreak{}.\allowbreak{}.\allowbreak{}.\allowbreak{})}
  with \code{For(\allowbreak{}Select)} and
  \code{To(\allowbreak{}\textquotedbl{}shop\_\allowbreak{}support\textquotedbl{})} &
  Restricted to SELECT and to members of the role \\
  \tablerule
  Maintenance per tenant & Background job that iterates the tenants and sets each one
  (chapter~\ref{ch:using-the-plugin-in-asp}) & Runs through the normal policies \\
\end{datatable}

\begin{sourcecode}[Reporting role (\allowbreak{}tests/\allowbreak{}.\allowbreak{}.\allowbreak{}.\allowbreak{}Integration\allowbreak{}Tests/\allowbreak{}Cross\allowbreak{}Tenant\allowbreak{}Access\allowbreak{}Tests.\allowbreak{}cs)]{sql}
-- Created by an administrator: BYPASSRLS requires a superuser
CREATE ROLE shop_reporting LOGIN PASSWORD '...' NOSUPERUSER BYPASSRLS;
GRANT USAGE ON SCHEMA public TO shop_reporting;
GRANT SELECT ON ALL TABLES IN SCHEMA public TO shop_reporting;
\end{sourcecode}

The application role, in contrast, cannot switch the policies off. \code{SET
row\_\allowbreak{}security = off} does not disable RLS for a role without \code{BYPASSRLS};
PostgreSQL instead rejects every query the policies would affect.
\code{Cross\allowbreak{}Tenant\allowbreak{}Access\allowbreak{}Tests} checks all three: the
reporting role sees all tenants, it cannot write, and the app role cannot help itself.

\begin{important}[Audit cross-tenant access]
  Every connection of the reporting role is an access to all tenants. Log it
  (\code{log\_\allowbreak{}connections}, pgaudit) and grant the role to as few people and
  services as possible. The plugin's \code{Create\allowbreak{}Application\allowbreak{}Role}
  deliberately always creates roles with \code{NOBYPASSRLS}: bypass roles do not belong in
  the application's migrations.
\end{important}

% ------------------------------------------------------------
\chapter{Tenant-safe relationships}\label{ch:tenant-safe-relationships}

Which table mappings the plugin protects and which it rejects while the model is built is set by
the matrix on page~\pageref{sec:supported-table-mappings}; this chapter does not repeat it. It is
about the second axis: relationships between protected tables. RLS governs which \emph{rows} a
session sees. Whether a row only refers to rows of the same tenant is not governed by RLS. That
is \emph{referential tenant consistency}, and it comes from the model.

The sample uses customer, order and item: \code{Customer} → \code{Order} →
\code{Order\allowbreak{}Item}, all three \code{ITenant\allowbreak{}Owned}, plus a shared
\code{Product} without a tenant.

\section{Visibility is not consistency}

PostgreSQL documents explicitly that checks of foreign keys, unique and primary keys bypass RLS
to maintain data integrity, and warns of “covert channels” through these checks
(\href{https://www.postgresql.org/docs/17/ddl-rowsecurity.html}{PostgreSQL: Row Security Policies}).
The SQL Server documentation says nothing explicit about it;
\code{Sql\allowbreak{}Server\allowbreak{}Threat\allowbreak{}Model\allowbreak{}Tests} shows that
the foreign key check there also sees rows the filter predicate hides. With foreign keys on the
ID alone, as EF creates them by default, this follows:

\begin{datatable}{H{0.70}H{1.30}}
  \textbf{Observation} & \textbf{Cause} \\
  \tablehead
  Tenant A adds an item to an order of B & The FK check finds the order, although A cannot
  see it. \\
  \tablerule
  A learns whether a foreign ID exists & Unknown IDs fail on the foreign key, foreign ones
  do not. \\
  \tablerule
  The own item vanishes from \code{Include(\allowbreak{}i => i.\allowbreak{}Order)} & A required
  navigation becomes an \code{INNER JOIN}, and the order is invisible. \\
  \tablerule
  A deletes an own order, an item of B disappears & The cascading delete is a referential action
  and bypasses RLS as well. \\
  \tablerule
  A learns that an e-mail address exists at B & A global unique index reports the
  collision. \\
\end{datatable}

\code{Relationship\allowbreak{}Tests} proves every row against PostgreSQL with a restricted
role. None of these cases is a read leak in the strict sense; A sees no foreign row. But the
data loses its tenant boundary, and the error messages answer questions A should not be able
to ask.

\section{Composite foreign keys}

The solution is to make the tenant part of every reference. The principal gets an alternate key
of tenant and ID, the dependent side a foreign key of its tenant column and the ID:

% source: tests/Acme.EntityFrameworkCore.Tenancy.IntegrationTests/RelationshipTests.cs
\begin{sourcecode}[tests/\allowbreak{}.\allowbreak{}.\allowbreak{}.\allowbreak{}Integration\allowbreak{}Tests/\allowbreak{}Relationship\allowbreak{}Tests.\allowbreak{}cs › Tenant\allowbreak{}Safe\allowbreak{}Context]{csharp}
modelBuilder.Entity<Customer>(b =>
{
    b.HasAlternateKey(c => new { c.TenantId, c.Id });
    b.HasIndex(c => new { c.TenantId, c.Email }).IsUnique();   // unique per tenant
});

modelBuilder.Entity<Order>(b =>
{
    b.HasAlternateKey(o => new { o.TenantId, o.Id });
    b.HasOne(o => o.Customer).WithMany(c => c.Orders)
        .HasForeignKey(o => new { o.TenantId, o.CustomerId })
        .HasPrincipalKey(c => new { c.TenantId, c.Id })
        .OnDelete(DeleteBehavior.Restrict);                   // customers are not cascaded
});

modelBuilder.Entity<OrderItem>()
    .HasOne(i => i.Order).WithMany(o => o.Items)
    .HasForeignKey(i => new { i.TenantId, i.OrderId })
    .HasPrincipalKey(o => new { o.TenantId, o.Id })
    .OnDelete(DeleteBehavior.Cascade);                         // items belong to the order
\end{sourcecode}

Now an item of A refers to the pair (A, ID). An order of B never satisfies it, and the error is
the same as for an unknown ID: there is nothing left to learn. The tenant column is part of the
foreign key and the tenant property at the same time. The SaveChanges interceptor sets it on new
items, and fixup through the navigation sets the same value. No index is needed for the policy,
because key and foreign key start with the tenant column; the index convention recognizes that.

The price is an additional unique index per principal and a wider foreign key. The plugin does
not create these keys automatically. A convention could, but it would have to rewrite every
relationship between protected types and reconcile it with explicit configuration. The sample
leaves the decision to the application's model.

\section{Uniqueness per tenant}

A unique index over a business column applies to the whole table, that is across all tenants.
Uniqueness per tenant is almost always what is meant: \code{(TenantId, Email)} instead of
\code{Email}. Where uniqueness really has to be global, for example for the tenants'
subdomains, the column belongs in the tenant directory, which deliberately holds no tenant data.

\section{Shared master data}

\code{Product} has no tenant, no policy and is readable by everyone. Foreign keys from tenant data
to master data are unproblematic; the test reads the product through the item. The opposite
direction, master data with references to tenant data, must not exist. Writing to master data
should not be allowed to the app role; \code{Grant\allowbreak{}Table\allowbreak{}Privileges}
grants privileges on all tables of a schema, a finer grant is up to the application's
migrations.

\section{Join tables}

The plugin rejects implicit many-to-many tables when one side is tenant-owned. A join entity with
a tenant property of its own is protected (page~\pageref{sec:supported-table-mappings}). For
consistency it gets the same composite foreign keys as every dependent table:
\code{(TenantId, LeftId)} and \code{(TenantId, RightId)} on the alternate keys of both sides.
Then a link can only connect two rows of the same tenant.

\section{Delete rules}

Cascading deletes run in the database as referential actions and bypass RLS. With composite
foreign keys they still stay within the tenant, because every dependent row carries the same
tenant. Without them a cascade deletes foreign rows; the test
\code{Naive\allowbreak{}\_cascade\allowbreak{}\_deletes\allowbreak{}\_rows\allowbreak{}\_of\allowbreak{}\_another\allowbreak{}\_tenant}
shows exactly that. For roots such as the customer, \code{Restrict} is usually right, for parts of
an aggregate such as items \code{Cascade}. EF also cascades in the change tracker, but only for
loaded entities; the rule in the database is what counts.


% ------------------------------------------------------------
\chapter{Retrofitting and operating multi-tenancy}\label{ch:retrofitting-and-operating}

The previous chapters assume a new database. Usually, however, there already are tables with
data, a running application and users who keep working during the change. This chapter shows the
way from a table without a tenant column to RLS, and after that operation under retries and lost
connections. All steps are executable in \code{Operations\allowbreak{}Tests} against PostgreSQL.

\section{From a table without a tenant to RLS}

Whoever simply lets an existing entity implement \code{ITenant\allowbreak{}Owned} gets a single
migration: a column with the empty tenant as default, RLS, policy, index. After it nobody sees the
old rows any more. A backfill in the same migration does not reliably help either: with
\code{FORCE ROW LEVEL SECURITY} the policy also applies to the migration role, and without a tenant
set its \code{UPDATE} finds no row. The way therefore leads through two migrations, following the
pattern expand, fill, protect:

\begin{datatable}{H{0.22}H{0.90}H{1.88}}
  \textbf{No.} & \textbf{Step} & \textbf{Approach} \\
  \tablehead
  1 & Expand & The entity implements \code{ITenant\allowbreak{}Owned} but stays excluded with
  \code{Has\allowbreak{}No\allowbreak{}Row\allowbreak{}Level\allowbreak{}Security(\allowbreak{})}.
  The migration only adds the column. \\
  \tablerule
  2 & Fill & As owner, in the same or a migration of its own: every row gets its tenant from
  whatever knows the assignment, for example a mapping table or the tenant of the parent row. \\
  \tablerule
  3 & Switch the application & An app version whose model no longer has the exclusion: the
  interceptor sets the tenant on new rows. The database does not protect yet; the model does not
  check that. \\
  \tablerule
  4 & Close the gap & Assign rows that were created without a tenant between steps 1 and 3. \\
  \tablerule
  5 & Protect & Second migration: exclusion removed, so \code{ENABLE}/\code{FORCE}, policy and
  index. First of all a guard that stops as long as a row has no tenant. \\
  \tablerule
  6 & Tighten & Optionally composite keys and unique indexes per tenant
  (chapter~\ref{ch:tenant-safe-relationships}). \\
\end{datatable}

The guard is a step of its own at the start of \code{Up}, before the operations EF generates:

\begin{sourcecode}[Guard (\allowbreak{}SQL)]{postgresql}
DO $$ BEGIN
  IF EXISTS (SELECT 1 FROM "Tickets"
             WHERE "TenantId" = '00000000-0000-0000-0000-000000000000') THEN
    RAISE EXCEPTION 'Tickets without a tenant: backfill first';
  END IF;
END $$;
\end{sourcecode}

\begin{sourcecode}[Migrations/\allowbreak{}…\allowbreak{}\_Enable\allowbreak{}Ticket\allowbreak{}Isolation.\allowbreak{}cs]{csharp}
protected override void Up(MigrationBuilder migrationBuilder)
{
    // Before anything else: do not protect without a complete assignment
    migrationBuilder.Sql(TicketGuardSql);      // the guard above, as a constant

    // ... generated by EF: ENABLE/FORCE ROW LEVEL SECURITY, CREATE POLICY, CREATE INDEX
}
\end{sourcecode}

The migration runs in a transaction. If the guard fires, the table stays without RLS, and no
half-done state arises;
\code{Guard\allowbreak{}\_stops\allowbreak{}\_enabling\allowbreak{}\_rls\allowbreak{}\_while\allowbreak{}\_rows\allowbreak{}\_lack\allowbreak{}\_a\allowbreak{}\_tenant}
checks this. After a complete assignment every tenant sees exactly its rows.

\subsection*{App versions, Down and rollback}

During the change two app versions often run at the same time, and each meets one of the two
states of the database:

\begin{datatable}{H{0.95}H{1.00}H{1.05}}
  \textbf{Database} & \textbf{App without tenant} & \textbf{App with tenant} \\
  \tablehead
  After migration 1 (column, no RLS) & Runs; new rows without a tenant (step 4) & Runs, sets
  the tenant; not yet protected \\
  \tablerule
  After migration 2 (RLS) & Fail closed: reads nothing, writes fail & Protected \\
\end{datatable}

The old app version must therefore no longer run once migration 2 is applied. Rolling the
application back to a version without a tenant thus first requires the Down of migration 2. The
Down removes policy and RLS; column and values stay, but the table is unprotected again. The test
checks both. A Down of migration 1 drops the column and with it the assignment; it belongs in
development only.

Policy changes follow the same pattern. If the setting name changes
(\code{Use\allowbreak{}Setting\allowbreak{}Name}), the migration replaces the policy in its
transaction. Between the migration and the new app version the policy reads a setting the running
app does not set, and the database refuses everything until both match. A transition in which the
interceptor sets both settings is conceivable but not implemented. For SQL Server,
page~\pageref{sec:sqlserver-session-settings} describes the same case.

\section{Retries, transactions and lost connections}

The plugin sets the tenant when EF opens a connection. This determines the behavior under EF's
resilience mechanisms:

\begin{itemize}
  \item \textbf{Retries.} An execution strategy such as
    \code{Enable\allowbreak{}Retry\allowbreak{}On\allowbreak{}Failure} repeats an operation after
    a transient error on a newly opened connection. The interceptor runs again, the tenant is set.
    The test ends the session on the server with \code{pg\_\allowbreak{}terminate\_\allowbreak{}backend};
    the next query fails, is repeated and returns only rows of the tenant.
  \item \textbf{Explicit transactions.} With a retrying strategy EF refuses commands in a
    transaction the application began itself, because a retry would have to repeat the whole
    transaction. The transaction therefore belongs into
    \code{Create\allowbreak{}Execution\allowbreak{}Strategy(\allowbreak{}).\allowbreak{}Execute\allowbreak{}Async(\allowbreak{}…)}.
    Inside, EF opens the connection once, sets the tenant and keeps it until the commit. If the
    connection breaks in the middle of the transaction, the strategy repeats the whole delegate on
    a new connection; side effects outside the database then have to be idempotent.
  \item \textbf{Session reinitialization.} Npgsql resets a connection's state when it returns to
    the pool, as long as \code{No Reset On Close} is not enabled
    (\href{https://www.npgsql.org/doc/connection-string-parameters.html}{Npgsql: Connection String Parameters}).
    SQL Server also clears the session context when the connection is returned. The plugin relies
    on neither, but sets the tenant on every open; only connections that EF does not open need
    \code{Apply\allowbreak{}Tenant\allowbreak{}Session} (chapter~\ref{ch:interceptors}).
\end{itemize}

Not implemented are transaction-local settings for transaction pooling (section “Transaction
pooling” in chapter~\ref{ch:row-level-security-in-production}) and a read-only session key
(section “The threat model”). \code{Operations\allowbreak{}Tests} contains the lost connection and
the transaction inside the strategy; the cases of open connections are in
\code{Connection\allowbreak{}Lifetime\allowbreak{}Tests}.


% ------------------------------------------------------------
\chapter{Plugin performance}\label{ch:plugin-performance}

Plugin code runs at three frequencies. Knowing which code runs how often tells you where
optimization pays off:

\begin{datatable}{H{0.75}H{1.25}H{1.00}}
  \textbf{Frequency} & \textbf{Code} & \textbf{Consequence} \\
  \tablehead
  Once per provider & Conventions, type mapping plugins, singleton options & Readability
  over micro-optimization \\
  \tablerule
  Once per query shape & Translators, query interceptors & Moderate; the result is cached
  (chapter~\ref{ch:the-query-pipeline}) \\
  \tablerule
  On every execution & Interceptors for commands, connections, SaveChanges, materialization;
  evaluatable filter when the funcletizer extracts the parameter values & Hot path: no reflection, no unnecessary allocations \\
\end{datatable}

This leads to a few rules the sample follows: \code{Method\allowbreak{}Info} and
\code{Property\allowbreak{}Info} lookups live in static fields. The SaveChanges interceptor
only collects entries when there are tenant-relevant changes. The connection interceptor
costs one round trip per open, which is documented. Such assumptions should be measured, not
guessed. The repository contains a BenchmarkDotNet project for the SaveChanges logic; it
runs without a database:

% source: benchmarks/Acme.EntityFrameworkCore.Tenancy.Benchmarks/SaveChangesInterceptorBenchmark.cs
\begin{sourcecode}[benchmarks/\allowbreak{}.\allowbreak{}.\allowbreak{}.\allowbreak{}/\allowbreak{}Save\allowbreak{}Changes\allowbreak{}Interceptor\allowbreak{}Benchmark.\allowbreak{}cs (\allowbreak{}shortened)]{csharp}
[MemoryDiagnoser]
public class SaveChangesInterceptorBenchmark
{
    [Params(1, 100)]
    public int Entities { get; set; }

    // GlobalSetup: context with UseAcmeTenancy() (no connection) and the interceptor

    [Benchmark(Baseline = true)]
    public int TrackOnly()                  // what EF does anyway before saving
    {
        AddOrders();
        var count = _db.ChangeTracker.Entries().Count();
        _db.ChangeTracker.Clear();
        return count;
    }

    [Benchmark]
    public int TrackAndApplyTenant()        // the same plus the plugin's SaveChanges logic
    {
        AddOrders();
        _interceptor.ApplyTenant(_db);
        var count = _db.ChangeTracker.Entries().Count();
        _db.ChangeTracker.Clear();
        return count;
    }
}
\end{sourcecode}

\begin{sourcecode}[Running the benchmark]{text}
dotnet run -c Release --project benchmarks/Acme.EntityFrameworkCore.Tenancy.Benchmarks
\end{sourcecode}

\begin{datatable}{H{1.67}H{0.83}H{0.83}H{0.83}H{0.83}}
  \textbf{Method} & \textbf{Entities} & \textbf{Mean} & \textbf{Ratio} & \textbf{Allocated} \\
  \tablehead
  \code{Track\allowbreak{}Only} & 1 & 1.5 µs & 1.00 & 1.8 KB \\
  \tablerule
  \code{Track\allowbreak{}And\allowbreak{}Apply\allowbreak{}Tenant} & 1 & 2.3 µs & 1.54 &
  2.7 KB \\
  \tablerule
  \code{Track\allowbreak{}Only} & 100 & 123 µs & 1.00 & 134 KB \\
  \tablerule
  \code{Track\allowbreak{}And\allowbreak{}Apply\allowbreak{}Tenant} & 100 & 170 µs & 1.38 &
  202 KB \\
\end{datatable}

Example results (short job, .NET 10, author's machine; absolute values vary, the ratios are
more meaningful). The plugin adds about half a microsecond per saved entity, roughly 40 to
50 percent on top of change tracking alone. That sounds like much, but a single round trip
to the database costs hundreds of microseconds; the connection interceptor's extra round
trip per open (chapter~\ref{ch:interceptors}) therefore weighs far more than the SaveChanges
logic. That is the kind of insight a benchmark provides: it shows where optimizing is
worthwhile and where it is not.

\begin{tip}
  For the startup of large models, the compiled model
  (chapter~\ref{ch:compiled-models-trimming-and-nativeaot}) is the most effective lever. A
  plugin's conventions then no longer run at run time at all; that is why it pays to test
  compiled model support rather than merely claim it.
\end{tip}

% ------------------------------------------------------------
\chapter{Using the plugin in ASP.NET Core}\label{ch:using-the-plugin-in-asp}

The sample \code{samples/\allowbreak{}Shop\allowbreak{}Api} shows the way from the plugin
into a web API. The tenant comes from the claim \code{tenant\_\allowbreak{}id} of the
signed-in user:

% source: samples/ShopApi/HttpTenantProvider.cs
\begin{sourcecode}[samples/\allowbreak{}Shop\allowbreak{}Api/\allowbreak{}Http\allowbreak{}Tenant\allowbreak{}Provider.\allowbreak{}cs]{csharp}
public sealed class HttpTenantProvider : ITenantProvider
{
    public const string ClaimType = "tenant_id";

    private readonly IHttpContextAccessor _httpContextAccessor;

    public HttpTenantProvider(IHttpContextAccessor httpContextAccessor)
    {
        ArgumentNullException.ThrowIfNull(httpContextAccessor);

        _httpContextAccessor = httpContextAccessor;
    }

    // Determined anew on every access: also works as a singleton (pooling)
    public TenantId? CurrentTenantId
    {
        get
        {
            var value = _httpContextAccessor.HttpContext?.User.FindFirst(ClaimType)?.Value;
            return Guid.TryParse(value, out var id) ? new TenantId(id) : null;
        }
    }
}
\end{sourcecode}

% source: samples/ShopApi/Program.cs
\begin{sourcecode}[samples/\allowbreak{}Shop\allowbreak{}Api/\allowbreak{}Program.\allowbreak{}cs (\allowbreak{}excerpt)]{csharp}
builder.Services.AddHttpContextAccessor();
builder.Services.AddAcmeTenancy<HttpTenantProvider>();
builder.Services.AddDbContext<ShopApiContext>((sp, options) => options
    .UseNpgsql(builder.Configuration.GetConnectionString("Shop"))
    .UseAcmeTenancy(sp));

// No tenant filter in code: RLS in the database provides the separation
app.MapGet("/orders", async (ShopApiContext db) =>
    await db.Orders.OrderBy(o => o.Id).ToListAsync());
\end{sourcecode}

\begin{important}[Where the tenant comes from]
  The tenant must come from a trustworthy source: a claim from a signed token. A header or
  query parameter that the client sets freely undermines the entire isolation. The sample
  contains a header simulation for the development environment only and labels it
  explicitly.
\end{important}

For \code{Add\allowbreak{}Db\allowbreak{}Context\allowbreak{}Pool}, register the provider as
a singleton. This is possible because \code{Http\allowbreak{}Tenant\allowbreak{}Provider}
determines the tenant anew on every access; the interceptors ask it on every connection and
every save:

\begin{sourcecode}[Variant with pooling]{csharp}
builder.Services.AddAcmeTenancy<HttpTenantProvider>(ServiceLifetime.Singleton);
builder.Services.AddDbContextPool<ShopApiContext>((sp, options) => options
    .UseNpgsql(connectionString)
    .UseAcmeTenancy(sp));   // sp: root provider; interceptors are created once
\end{sourcecode}

\section{DbContext factories}

\code{Add\allowbreak{}Db\allowbreak{}Context\allowbreak{}Factory} is the usual choice in
Blazor Server and for parallel work. Unlike \code{Add\allowbreak{}Db\allowbreak{}Context},
it registers the factory and the options as \textbf{singletons} by default. The options
callback then receives the root provider, not the scope of the current request. For a scoped
tenant provider that holds state per scope, this has consequences:

\begin{datatable}{H{0.87}H{1.13}}
  \textbf{Registration} & \textbf{Behavior} \\
  \tablehead
  Singleton factory, scope validation on (ASP.NET Core in Development) & Exception: “Cannot
  resolve scoped service 'ITenantProvider' from root provider.” \\
  \tablerule
  Singleton factory, scope validation off (production default) & No error. The interceptors
  get an instance from the root container that no request ever fills: no tenant, saving
  throws, queries see nothing \\
  \tablerule
  Factory registered as \code{Scoped} & Every scope gets its own factory with its own tenant
  provider: correct \\
\end{datatable}

\begin{sourcecode}[Factory with the tenant of the request]{csharp}
builder.Services.AddDbContextFactory<ShopApiContext>(
    (sp, options) => options.UseNpgsql(connectionString).UseAcmeTenancy(sp),
    ServiceLifetime.Scoped);         // factory and options per scope
\end{sourcecode}

A provider that determines the tenant anew on every access, such as
\code{Http\allowbreak{}Tenant\allowbreak{}Provider}, works in every variant; the problem
affects providers with state per scope.
\code{Add\allowbreak{}Pooled\allowbreak{}Db\allowbreak{}Context\allowbreak{}Factory} always
behaves like pooling: the options exist once, so the tenant provider must be such a
singleton. \code{Db\allowbreak{}Context\allowbreak{}Factory\allowbreak{}Tests} shows all
three variants.

\section{Outside of HTTP requests}

Background services, message handlers and scheduled jobs have no request and no claim. The
tenant then comes from the job itself, for example from the message or from a list of
tenants to process. The sample's \code{Ambient\allowbreak{}Tenant\allowbreak{}Provider}
stores it in an \code{Async\allowbreak{}Local}, which flows with the asynchronous control
flow:

% source: samples/ShopSample/AmbientTenantProvider.cs
\begin{sourcecode}[samples/\allowbreak{}Shop\allowbreak{}Sample/\allowbreak{}Ambient\allowbreak{}Tenant\allowbreak{}Provider.\allowbreak{}cs (\allowbreak{}shortened)]{csharp}
public sealed class AmbientTenantProvider : ITenantProvider
{
    private static readonly AsyncLocal<TenantId?> Current = new();

    public TenantId? CurrentTenantId => Current.Value;

    public static IDisposable Use(TenantId tenant)
    {
        var previous = Current.Value;
        Current.Value = tenant;
        return new Restore(previous);     // sets the previous tenant back on Dispose
    }

    // private sealed class Restore : IDisposable { ... Current.Value = _previous; }
}
\end{sourcecode}

\begin{sourcecode}[Job pattern (\allowbreak{}tests/\allowbreak{}.\allowbreak{}.\allowbreak{}.\allowbreak{}/\allowbreak{}Background\allowbreak{}Tenant\allowbreak{}Tests.\allowbreak{}cs)]{csharp}
// Tenant first, then scope and context: the interceptors are created inside the scope
using (AmbientTenantProvider.Use(tenant))
{
    await using var scope = scopeFactory.CreateAsyncScope();
    var db = scope.ServiceProvider.GetRequiredService<ShopContext>();
    // ... work for this tenant ...
}
\end{sourcecode}

Because the provider reads the \code{Async\allowbreak{}Local} on every access, it can be
registered as a singleton. Parallel jobs keep their own tenant, and after the \code{using}
block no tenant is left behind; \code{Background\allowbreak{}Tenant\allowbreak{}Tests}
checks both with 20 parallel jobs. An application with web requests \emph{and} jobs combines
both sources in one provider: the claim first, otherwise the ambient tenant.

\begin{important}[Never fall back to “all tenants”]
  Without a tenant, a job must fail or see nothing, never everything. The plugin behaves
  this way: the session value stays empty, the policies match no rows and saving throws.
  Jobs that really need all tenants use the reporting role from
  chapter~\ref{ch:row-level-security-in-production}.
\end{important}

% ------------------------------------------------------------
\chapter{CI and API stability}\label{ch:ci-and-api-stability}

\section{The workflow}

The GitHub Actions workflow in the repository implements the recommendations from
chapters~\ref{ch:test-strategy} and \ref{ch:versioning-and-compatibility}. One job builds
and tests without and with a database, a second builds the handbook:

\begin{sourcecode}[.\allowbreak{}github/\allowbreak{}workflows/\allowbreak{}ci.\allowbreak{}yml (\allowbreak{}shortened)]{text}
jobs:
  build:
    runs-on: ubuntu-latest
    steps:
      - uses: actions/checkout@v4
      - uses: actions/setup-dotnet@v4
        with: { global-json-file: global.json }
      - run: dotnet build -c Release
      - run: dotnet test tests/Acme.EntityFrameworkCore.Tenancy.Tests -c Release --no-build
      - run: dotnet test tests/Acme.EntityFrameworkCore.Tenancy.IntegrationTests
          -c Release --no-build
      - run: dotnet test tests/Acme.EntityFrameworkCore.Tenancy.SqlServer.Tests
          -c Release --no-build

  handbook:                            # both PDFs with LuaLaTeX, in docs/handbook
    runs-on: ubuntu-latest
    container: texlive/texlive@sha256:540d…   # TeX Live 2026, pinned by digest
    steps:
      # ... checkout ...
      - run: python3 check_code.py       # line length, excerpts against the sources
      - run: python3 check_structure.py  # editions in sync, references, tables
      - run: latexmk EF-Core-Plugin-Handbook.tex

  consumer:                            # check the packages as a consumer
    steps:
      - run: eng/consumer-test/run.sh --nativeaot
\end{sourcecode}

Code excerpts go stale silently when the code changes. A comment line
\code{\% source: <path>} directly before the listing therefore binds an excerpt to its source
file, and \code{check\_\allowbreak{}code.\allowbreak{}py} compares the two: comments and elisions
(\code{...}) do not count, so the German edition may translate comments and excerpts may be
shortened; every other line must appear in the source in the same order, line breaks may
differ. With the suffix \code{(strings)} translated messages are left out as well. An excerpt
whose title names a file of the repository and that is not bound must be called “simplified”.
\code{check\_\allowbreak{}structure.\allowbreak{}py} checks that both editions have the same
sections, listings, bindings, tables and callouts per chapter, that every reference target exists
and that the column widths of every table add up.

The sample builds against stable versions only (.NET 10, EF Core 10). A second workflow,
\code{compatibility.\allowbreak{}yml}, runs weekly and tests against the newest stable 10.0.x
patches, because the regular CI builds against the beginning of the version range.
Chapter~\ref{ch:versioning-and-compatibility} sketches a channel for the next EF major version with previews and
RCs; it needs the matching .NET prerelease SDK, because the prereleases of EF Core 11 only
offer \code{net11.\allowbreak{}0}. That is a property of this version, not a general rule:
EF Core 9 targeted .NET 8.

\section{Protecting the public API}

A plugin has two public APIs: the one for users (Fluent API, options, types) and the one for
extenders (virtual methods, dependencies, public services). Both break easily and unnoticed.
The
\code{Microsoft.\allowbreak{}Code\allowbreak{}Analysis.\allowbreak{}Public\allowbreak{}Api\allowbreak{}Analyzers}
package keeps the public API in two text files
(\code{Public\allowbreak{}API.\allowbreak{}Shipped.\allowbreak{}txt},
\code{Public\allowbreak{}API.\allowbreak{}Unshipped.\allowbreak{}txt}) and reports every
deviation. Every API change thus becomes visible as a diff in the pull request.

\begin{note}[Setup]
  On the first build, the analyzer reports every public API as RS0016 because the files are
  empty. The code fix “Add all items … to the public API” fills them, on the command line with
  \code{dotnet format analyzers -{}-diagnostics RS0016 -{}-severity warn}. At release time,
  the content moves from Unshipped to Shipped; from then on, removing a shipped API is a
  visible breaking change.
\end{note}

\begin{important}[Warnings are not enough]
  As long as RS0016 and RS0017 are only warnings, they get lost in the build output. In the
  sample, the files thus stayed empty up to edition~1.0: 178 warnings in a full build, none
  in incremental ones, and nobody noticed.
  \code{Directory.\allowbreak{}Build.\allowbreak{}props} therefore turns both into errors
  (\code{Warnings\allowbreak{}As\allowbreak{}Errors}); a new public API now has to be entered
  in \code{Public\allowbreak{}API.\allowbreak{}Unshipped.\allowbreak{}txt}, otherwise the
  build breaks.
\end{important}
